\begin{longtable}{@{}L{0.9cm}L{4.3cm}L{10cm}@{}}
\toprule
\textbf{No.} & \textbf{Decision} & \textbf{Alternatives and reasoning} \\
\midrule\endhead
D1 & Java 21 + Spring Boot 3 & Alternatives: Node/Express, Django, plain Java Swing. Why: course is OOP in Java; Spring's DI makes Strategy/Observer/Repository natural and swappable; mature security + JPA; industry-relevant. \\
D2 & Server-rendered Thymeleaf UI (not React SPA) & Alternatives: React + REST. Why: one codebase/language for a short timeline; role-specific views easy with Spring Security; REST endpoints can still be added for intelligent features. Limitation: less dynamic UI. \\
D3 & PostgreSQL (prod/demo) + H2 (tests) & Alternatives: MySQL, MongoDB. Why: relational data (users, incidents, SLAs, audit) with strong integrity; JPA keeps the code DB-agnostic so tests run fast in-memory. \\
D4 & SINGLE\_TABLE inheritance for the User hierarchy & Alternatives: JOINED, TABLE\_PER\_CLASS, a single User class with a role field. Why: Employee/Technician/Administrator are genuine is-a subtypes with only a few extra columns each. One table makes login a single query with no joins and keeps the JPA mapping simple. Trade-off: subtype-only columns must be nullable. \\
D5 & Stateless State objects mapped from a persisted status enum & Alternatives: storing state objects directly, a big switch on status in Incident. Why: JPA persists the simple \code{IncidentStatus} enum, and \code{IncidentStateRegistry} (an \code{EnumMap}) maps it to one shared, stateless \code{IncidentState} object. Each state overrides only the transitions it allows, so adding a state touches one class plus the registry, and illegal transitions fail in one place (\code{AbstractIncidentState}). \\
D6 & A shared \code{BaseEntity} superclass for identity & Alternatives: repeat \code{id}, \code{version}, \code{equals}/\code{hashCode} in every entity. Why: every entity genuinely \textit{is a} persistent, identifiable object, so this is an is-a relationship, not inheritance just to reuse code. One place defines id-based equality (safe for Hibernate proxies) and optimistic locking (\code{@Version}), so two admins editing the same incident can't silently overwrite each other. It's a \code{@MappedSuperclass}, so there's no extra table. \\
D7 & Subtype-only user fields use wrapper types (\code{Integer}, \code{Boolean}) & Because of SINGLE\_TABLE (D4), a Technician's columns are NULL for Employee rows. Wrapper types accept NULL, and the getters translate NULL into safe defaults, so callers never see nulls. \\
D8 & Value objects for incident number, SLA record, resolution and feedback & \code{IncidentNumber} is a record (\code{INC1001}), so user-facing numbers never expose database ids. \code{SlaRecord}, \code{Resolution} and \code{Feedback} are \code{@Embeddable}: they're stored in the incident row and have no identity of their own, which is exactly composition in UML. \\
D9 & Who decides what in the incident lifecycle & \textbf{State objects} decide \textit{whether} an operation is legal in the current status and \textit{which} status comes next. They hold no data and never touch the incident's fields. **\code{Incident}\textbf{ asks its state first, then updates its own fields. If the state rejects the operation, nothing has been changed yet. }\code{IncidentService}\textit{* decides }who* may do it (reporter, assigned technician, administrator) and publishes an event afterwards. Alternatives: one big \code{switch} on status inside \code{Incident}, or states that change the... \\
D10 & Domain events via Spring's \code{ApplicationEventPublisher} (Observer) & Alternatives: a hand-written listener list inside \code{IncidentService}, or calling audit and notification code directly. Why: Spring already provides a well-tested publish/subscribe mechanism. \code{IncidentService} only publishes an immutable \code{IncidentEvent} and doesn't know who listens, so audit (M6), notifications (M6) and escalation (M8) can be added without changing it. \\
D11 & Priority as a Strategy, selected by configuration & Alternatives: one hard-coded formula inside \code{IncidentService}. Why: the brief explicitly allows teams to choose their own algorithm. \code{WeightedScoringPriorityStrategy} (default) and \code{MatrixPriorityStrategy} (ITIL matrix) share one interface, and \code{PriorityService} picks the one named in \code{itsm.priority.strategy}. A new algorithm, such as one based on historical data, is one new \code{@Component}. Every result carries an explanation string, so users can see why an incident is P1. \\
D12 & SLA rules live in \code{SlaRecord}; \code{Incident} drives them & Alternatives: a separate service that edits SLA fields from outside. Why: SLA timing is part of the incident's own state. \code{Incident.acknowledge()} records the response time and \code{Incident.resolve()} records the resolution time automatically, so no caller can forget. \code{SlaService} only chooses \textit{which} policy applies (the service's own, else the default). The rules: meeting a deadline exactly on time is not a breach, a breach is permanent, and it is reported only once. Known limitation: the SLA clock keeps running... \\
D13 & Classification happens automatically when the reporter picks a category & If the employee chooses a category when reporting, the system classifies the incident at once (priority calculated, SLA started), and the audit trail records the actor as \code{system}. Otherwise it stays NEW until an administrator picks a category. Either way the priority is always calculated, never typed in by hand. \\
D14 & Assignment strategies switchable at runtime, stored as a setting & Alternatives: a config-file choice (as for priority), or an \code{if/else} per strategy in the service. Why: the brief asks that the assignment strategy can be changed without redesigning the application, and an administrator should be able to do it from the UI. \code{AssignmentStrategyRegistry} receives every \code{AssignmentStrategy} from Spring and stores the active name in the \code{system\_setting} table, so the choice takes effect immediately and survives restarts. The four strategies are skill, load, round-robin and smart.... \\
D15 & Build output kept out of iCloud (\code{target.nosync}) & The project sits in \code{\textasciitilde{}/Documents}, which is synced by iCloud Drive. iCloud created duplicate files such as \code{UserFactory 2.class} in \code{target/} and offloaded others mid-build, which broke the tests with "Operation timed out". Maven now builds into \code{target.nosync/}, and iCloud does not sync folders whose names end in \code{.nosync}. Moving the whole project outside \code{\textasciitilde{}/Documents} would be safer still. \\
D16 & Audit and notifications are Observers, run in the same transaction & \code{AuditListener} and \code{NotificationListener} are Spring \code{@EventListener}s on \code{IncidentEvent}. They run synchronously, inside the transaction that changed the incident. \textbf{Audit:} an incident change and its audit line commit together or not at all, so the trail can never miss a change. \code{@Order} makes audit run before notifications. \textbf{Notifications:} every channel call is wrapped in try/catch, so a failing email or SMS provider is logged and skipped. It never rolls back the incident change or blocks the other... \\
D17 & Notification design: Strategy (channels) + Adapter (gateways) + Factory (messages) + rule table (recipients) & **\code{NotificationChannel}** (Strategy): in-app, email and SMS share one interface. Each channel decides in \code{accepts()} whether it applies; for example, SMS is used only for urgent messages to users with a phone number. **\code{EmailGateway} / \code{SmsGateway}** (Adapter targets): the channels depend only on these interfaces. The \code{Logging...Gateway} classes simulate a provider and keep an outbox for administrators. A real provider means one new adapter class. **\code{NotificationFactory}** (Factory): builds the message text per... \\
D18 & Web layer: three lines of security, thin controllers, open-session-in-view & \textbf{Security in depth:} 1. URL rules per role in \code{SecurityConfig} (\code{/employee/**}, \code{/tech/**}, \code{/admin/**}). 2. \code{@PreAuthorize("hasRole('ADMIN')")} on administration service methods. 3. Ownership checks in \code{IncidentService} (own incident, assigned incident). Passwords use BCrypt, CSRF protection is on for every form, Bean Validation guards all input, and the H2 console is reachable only in the \code{dev} profile. Invalid logins go to \code{/login?error}; inactive accounts cannot log in. \textbf{Thin controllers:} each action... \\
D19 & Search with JPA Specifications & Each filter is a small \code{Specification\textless{}Incident\textgreater{}} in \code{IncidentSpecifications}. It returns null when its form field is empty, so \code{Specification.allOf(...)} combines only the filters actually used. The search supports 10 criteria (number, text, from, to, status, priority, category, service, technician, SLA status) plus a visibility scope per role. Alternatives: one repository method per combination, which explodes in number, or hand-built JPQL strings, which are error-prone and open to injection. \\
D20 & Intelligent features A, B, C and E, all behind interfaces & \textbf{A, duplicates:} \code{TextNormalizer} handles stopwords, simple plurals and synonyms (transaction/checkout/pay $\rightarrow$ payment; down/failing/unable $\rightarrow$ outage). \code{TextSimilarityDuplicateDetector} scores Jaccard word overlap, +0.2 when the service is the same, and flags a match at 0.5 or above. A match always points at the \textit{original} report, not at another copy. The same employee repeating the same report within an hour gets \code{DuplicateIncidentException}. A person confirms or dismisses each suggestion; the system never... \\
D21 & Demo data drives the real domain model & \code{HistoricalIncidentSeeder} creates its 28 past incidents through the same constructors, \code{PriorityService}, \code{SlaService} and State-pattern methods as live incidents, just with past timestamps. Seeded data therefore always follows the business rules, and the dashboards, search and feature G work from the very first start. The history is numbered below INC1001 and is at least 2 days old, so it never interferes with the live demo (duplicate detection only looks back 24 h). For the video, a Payment Gateway P1... \\
D22 & Chart.js visual analytics on the Admin Dashboard & Alternatives: Server-rendered SVG charts, text-only tables, external Grafana dashboard. Why: Chart.js runs client-side with a single lightweight script (pinned to 4.4.1 and served from a WebJar, see D26) and renders beautiful responsive HTML5 canvases for priority breakdown (doughnut chart) and service distribution (bar chart). It reads the data already prepared by \code{ReportService} and directly fulfills the "Real-time operational dashboard, Analytics and visualization" innovation challenge without backend overhead. \\
D23 & Service dependency graph via self-referencing @ManyToMany & Alternatives: Flat services list, separate dependency table with string foreign keys. Why: Demonstrates OOP aggregation and graph traversal. \code{ITService} defines a self-referencing \code{@ManyToMany} set \code{dependencies} mapped to \code{service\_dependencies}. Calling \code{findByDependenciesContaining()} navigates the graph in reverse to find downstream services impacted by an outage, giving technicians and admins instant impact analysis when investigating incidents. \\
D24 & RFC-4180 compliant CSV export for audit logs and search results & Alternatives: Third-party Apache POI / OpenCSV library. Why: Built with zero external dependencies using Java string builders and HTTP streaming headers (\code{Content-Disposition: attachment}). Escapes quotes and commas according to RFC-4180. Lets administrators export historical audit compliance logs and custom incident query results directly to spreadsheets. \\
D25 & External AI Incident Advisor with Google Gemini API \& local heuristic fallback & Alternatives: Hard-coded advice text or external Python AI microservice. Why: Directly satisfies the evaluation rubric's requirement for "use of APIs (Google API, any other), integration with third-party services". Uses Java 21's native \code{HttpClient} with non-blocking timeout handling. When \code{GEMINI\_API\_KEY} is provided, it calls Google's \code{generateContent} REST endpoint for automated triage and root-cause hypothesis; when offline or unconfigured, it seamlessly falls back to a deterministic rule-based expert... \\
D26 & Dark mode with CSS colour tokens; front-end libraries as WebJars & Alternatives: a second stylesheet for dark mode; keep loading Bootstrap and Chart.js from a CDN. Why: Bootstrap 5.3 switches its own components with \code{data-bs-theme}, so our \code{app.css} only defines colour tokens once per theme and every badge, chart and timeline uses the tokens. The choice (light / dark / system) is stored in the browser and applied by a tiny inline script before the first paint, so there is no white flash. WebJars put the libraries in \code{pom.xml} with exact versions and serve them from the app, so... \\
\bottomrule
\end{longtable}
